% ---------------------------------------------------------------------------
% Discussion & Limitations (whole project). Include with % ---------------------------------------------------------------------------
% Discussion & Limitations (whole project). Include with % ---------------------------------------------------------------------------
% Discussion & Limitations (whole project). Include with % ---------------------------------------------------------------------------
% Discussion & Limitations (whole project). Include with \input{discussion_limitations}.
% No extra packages needed. About 3,900 characters without spaces.
% If this section goes in, delete the "Discussion" subsection at the end of
% graph_clustering.tex: it says the same things and saves ~1,000 characters.
% ---------------------------------------------------------------------------

\section{Discussion and Limitations}
\label{sec:discussion}

We set out to answer two questions: can the structure of who reads what surface
books that purchase overlap neglects, and can a book's text stand in for that
structure when a book has no readers yet? The answer to both is a qualified yes.

\paragraph{Long-tail recommendation.}
The graph recommenders find long-tail books three times as often as
``readers also bought'' (recall@10 of 0.027 against 0.009) and recommend 82\% of
the catalogue instead of 47\%. The reason lies in how the graph is built.
Jaccard weighting divides out the popularity of both books, and the disparity
filter judges every edge against its own book rather than against a global
threshold, so a niche title keeps its strongest links instead of losing all of
them. Raw co-occurrence counts, by contrast, reproduce the bestseller bias of
purchase overlap. The price is paid at the head, where recall falls from 0.072
to 0.044; blending the two signals recovers most of it and still raises overall
recall by 16\%.

\paragraph{What clustering adds.}
On the reader graph alone, our most important finding is also the least
flattering one. Restricting candidates to a reader's communities gives exactly
the same accuracy as recommending from the same backbone without communities,
because 91\% of the backbone's edge weight already lies inside them: Louvain
recovers structure the backbone already encodes. This changes once the
description similarity is added to the graph. Description links are largely new
(only 6\% of them are also reader links), so in the combined graph only 68\% of
the weight falls inside communities, and the restriction becomes a filter. The
community recommender then beats the same graph without communities by a small
but significant 1.3\%, and finds 6\% more long-tail books than reader-only
communities. The two signals complement rather than replace each other:
descriptions alone reach only a third of the reader graph's recall. Beyond
accuracy, the partition is strong (modularity 0.84), agrees with Amazon's genres
far more than chance (NMI 0.29 against 0.06), and reveals reader groups that cut
across genres, such as post-apocalyptic thrillers read alongside prepper
manuals. Personalisation also matters more than membership: ranking a
community's books by centrality, as first planned, did worse than recommending
the most popular books.

\paragraph{New books.}
For books without readers, SBERT similarity combined with popularity gives the
best hit rate (0.482 against 0.458 for popularity alone). The margin grows to
69\% on the strict test that requires at least two shared readers, which suggests
that descriptions capture shared taste and not merely popularity. The gain is
largest for the least-read new books, exactly the ones popularity cannot help,
but descriptions shorter than about 200 characters carry too little information
to be useful. TF-IDF, which only matches shared words, falls well behind SBERT.

\paragraph{Locality-sensitive hashing.}
LSH proved useful as a candidate generator but not as a final answer. MinHash
on titles proposed 10,728 duplicate pairs, yet only 819 survived a check on
title and description, so its candidates had to be verified. Random-hyperplane
LSH cut the number of comparisons by a factor of 6 to 60 while keeping up to
72\% of the exact nearest neighbours, but with 150k candidates an exact search
took only 18.5 seconds as one matrix multiplication. At this scale exact search
is the better choice; LSH would only pay off on a much larger catalogue.

\paragraph{Limitations.}
The data comes from the 2014 SNAP release and was drawn from its 5-core subset,
so every book had at least five reviews before sampling. Our long tail is
therefore the tail of an already filtered catalogue, and the results may
overstate how well any method handles truly obscure books. Reviews are only a
proxy for purchases, and a liked book is defined purely by a rating of at least
four. This is a coarse signal, since 81\% of all ratings are four or five stars,
and the review text itself (143 words on average) was never used to judge
preference.
% If the sentiment part is added, say here how it refines "liked".
The metadata has no author field and assigns a genre to only about one book in
five, so genre agreement is measured on a subset and series were detected
heuristically from titles. Only 70\% of the books have both a title and a
description, so the description graph leaves the rest isolated. Duplicate editions were merged only in the cold-start
part; the graph results treat editions as separate books, which splits their
readers. Finally, every evaluation is offline and uses a single time split
(January to July 2014), and absolute recall stays low, as expected when choosing
ten books out of 48,000.
.
% No extra packages needed. About 3,900 characters without spaces.
% If this section goes in, delete the "Discussion" subsection at the end of
% graph_clustering.tex: it says the same things and saves ~1,000 characters.
% ---------------------------------------------------------------------------

\section{Discussion and Limitations}
\label{sec:discussion}

We set out to answer two questions: can the structure of who reads what surface
books that purchase overlap neglects, and can a book's text stand in for that
structure when a book has no readers yet? The answer to both is a qualified yes.

\paragraph{Long-tail recommendation.}
The graph recommenders find long-tail books three times as often as
``readers also bought'' (recall@10 of 0.027 against 0.009) and recommend 82\% of
the catalogue instead of 47\%. The reason lies in how the graph is built.
Jaccard weighting divides out the popularity of both books, and the disparity
filter judges every edge against its own book rather than against a global
threshold, so a niche title keeps its strongest links instead of losing all of
them. Raw co-occurrence counts, by contrast, reproduce the bestseller bias of
purchase overlap. The price is paid at the head, where recall falls from 0.072
to 0.044; blending the two signals recovers most of it and still raises overall
recall by 16\%.

\paragraph{What clustering adds.}
On the reader graph alone, our most important finding is also the least
flattering one. Restricting candidates to a reader's communities gives exactly
the same accuracy as recommending from the same backbone without communities,
because 91\% of the backbone's edge weight already lies inside them: Louvain
recovers structure the backbone already encodes. This changes once the
description similarity is added to the graph. Description links are largely new
(only 6\% of them are also reader links), so in the combined graph only 68\% of
the weight falls inside communities, and the restriction becomes a filter. The
community recommender then beats the same graph without communities by a small
but significant 1.3\%, and finds 6\% more long-tail books than reader-only
communities. The two signals complement rather than replace each other:
descriptions alone reach only a third of the reader graph's recall. Beyond
accuracy, the partition is strong (modularity 0.84), agrees with Amazon's genres
far more than chance (NMI 0.29 against 0.06), and reveals reader groups that cut
across genres, such as post-apocalyptic thrillers read alongside prepper
manuals. Personalisation also matters more than membership: ranking a
community's books by centrality, as first planned, did worse than recommending
the most popular books.

\paragraph{New books.}
For books without readers, SBERT similarity combined with popularity gives the
best hit rate (0.482 against 0.458 for popularity alone). The margin grows to
69\% on the strict test that requires at least two shared readers, which suggests
that descriptions capture shared taste and not merely popularity. The gain is
largest for the least-read new books, exactly the ones popularity cannot help,
but descriptions shorter than about 200 characters carry too little information
to be useful. TF-IDF, which only matches shared words, falls well behind SBERT.

\paragraph{Locality-sensitive hashing.}
LSH proved useful as a candidate generator but not as a final answer. MinHash
on titles proposed 10,728 duplicate pairs, yet only 819 survived a check on
title and description, so its candidates had to be verified. Random-hyperplane
LSH cut the number of comparisons by a factor of 6 to 60 while keeping up to
72\% of the exact nearest neighbours, but with 150k candidates an exact search
took only 18.5 seconds as one matrix multiplication. At this scale exact search
is the better choice; LSH would only pay off on a much larger catalogue.

\paragraph{Limitations.}
The data comes from the 2014 SNAP release and was drawn from its 5-core subset,
so every book had at least five reviews before sampling. Our long tail is
therefore the tail of an already filtered catalogue, and the results may
overstate how well any method handles truly obscure books. Reviews are only a
proxy for purchases, and a liked book is defined purely by a rating of at least
four. This is a coarse signal, since 81\% of all ratings are four or five stars,
and the review text itself (143 words on average) was never used to judge
preference.
% If the sentiment part is added, say here how it refines "liked".
The metadata has no author field and assigns a genre to only about one book in
five, so genre agreement is measured on a subset and series were detected
heuristically from titles. Only 70\% of the books have both a title and a
description, so the description graph leaves the rest isolated. Duplicate editions were merged only in the cold-start
part; the graph results treat editions as separate books, which splits their
readers. Finally, every evaluation is offline and uses a single time split
(January to July 2014), and absolute recall stays low, as expected when choosing
ten books out of 48,000.
.
% No extra packages needed. About 3,900 characters without spaces.
% If this section goes in, delete the "Discussion" subsection at the end of
% graph_clustering.tex: it says the same things and saves ~1,000 characters.
% ---------------------------------------------------------------------------

\section{Discussion and Limitations}
\label{sec:discussion}

We set out to answer two questions: can the structure of who reads what surface
books that purchase overlap neglects, and can a book's text stand in for that
structure when a book has no readers yet? The answer to both is a qualified yes.

\paragraph{Long-tail recommendation.}
The graph recommenders find long-tail books three times as often as
``readers also bought'' (recall@10 of 0.027 against 0.009) and recommend 82\% of
the catalogue instead of 47\%. The reason lies in how the graph is built.
Jaccard weighting divides out the popularity of both books, and the disparity
filter judges every edge against its own book rather than against a global
threshold, so a niche title keeps its strongest links instead of losing all of
them. Raw co-occurrence counts, by contrast, reproduce the bestseller bias of
purchase overlap. The price is paid at the head, where recall falls from 0.072
to 0.044; blending the two signals recovers most of it and still raises overall
recall by 16\%.

\paragraph{What clustering adds.}
On the reader graph alone, our most important finding is also the least
flattering one. Restricting candidates to a reader's communities gives exactly
the same accuracy as recommending from the same backbone without communities,
because 91\% of the backbone's edge weight already lies inside them: Louvain
recovers structure the backbone already encodes. This changes once the
description similarity is added to the graph. Description links are largely new
(only 6\% of them are also reader links), so in the combined graph only 68\% of
the weight falls inside communities, and the restriction becomes a filter. The
community recommender then beats the same graph without communities by a small
but significant 1.3\%, and finds 6\% more long-tail books than reader-only
communities. The two signals complement rather than replace each other:
descriptions alone reach only a third of the reader graph's recall. Beyond
accuracy, the partition is strong (modularity 0.84), agrees with Amazon's genres
far more than chance (NMI 0.29 against 0.06), and reveals reader groups that cut
across genres, such as post-apocalyptic thrillers read alongside prepper
manuals. Personalisation also matters more than membership: ranking a
community's books by centrality, as first planned, did worse than recommending
the most popular books.

\paragraph{New books.}
For books without readers, SBERT similarity combined with popularity gives the
best hit rate (0.482 against 0.458 for popularity alone). The margin grows to
69\% on the strict test that requires at least two shared readers, which suggests
that descriptions capture shared taste and not merely popularity. The gain is
largest for the least-read new books, exactly the ones popularity cannot help,
but descriptions shorter than about 200 characters carry too little information
to be useful. TF-IDF, which only matches shared words, falls well behind SBERT.

\paragraph{Locality-sensitive hashing.}
LSH proved useful as a candidate generator but not as a final answer. MinHash
on titles proposed 10,728 duplicate pairs, yet only 819 survived a check on
title and description, so its candidates had to be verified. Random-hyperplane
LSH cut the number of comparisons by a factor of 6 to 60 while keeping up to
72\% of the exact nearest neighbours, but with 150k candidates an exact search
took only 18.5 seconds as one matrix multiplication. At this scale exact search
is the better choice; LSH would only pay off on a much larger catalogue.

\paragraph{Limitations.}
The data comes from the 2014 SNAP release and was drawn from its 5-core subset,
so every book had at least five reviews before sampling. Our long tail is
therefore the tail of an already filtered catalogue, and the results may
overstate how well any method handles truly obscure books. Reviews are only a
proxy for purchases, and a liked book is defined purely by a rating of at least
four. This is a coarse signal, since 81\% of all ratings are four or five stars,
and the review text itself (143 words on average) was never used to judge
preference.
% If the sentiment part is added, say here how it refines "liked".
The metadata has no author field and assigns a genre to only about one book in
five, so genre agreement is measured on a subset and series were detected
heuristically from titles. Only 70\% of the books have both a title and a
description, so the description graph leaves the rest isolated. Duplicate editions were merged only in the cold-start
part; the graph results treat editions as separate books, which splits their
readers. Finally, every evaluation is offline and uses a single time split
(January to July 2014), and absolute recall stays low, as expected when choosing
ten books out of 48,000.
.
% No extra packages needed. About 3,900 characters without spaces.
% If this section goes in, delete the "Discussion" subsection at the end of
% graph_clustering.tex: it says the same things and saves ~1,000 characters.
% ---------------------------------------------------------------------------

\section{Discussion and Limitations}
\label{sec:discussion}

We set out to answer two questions: can the structure of who reads what surface
books that purchase overlap neglects, and can a book's text stand in for that
structure when a book has no readers yet? The answer to both is a qualified yes.

\paragraph{Long-tail recommendation.}
The graph recommenders find long-tail books three times as often as
``readers also bought'' (recall@10 of 0.027 against 0.009) and recommend 82\% of
the catalogue instead of 47\%. The reason lies in how the graph is built.
Jaccard weighting divides out the popularity of both books, and the disparity
filter judges every edge against its own book rather than against a global
threshold, so a niche title keeps its strongest links instead of losing all of
them. Raw co-occurrence counts, by contrast, reproduce the bestseller bias of
purchase overlap. The price is paid at the head, where recall falls from 0.072
to 0.044; blending the two signals recovers most of it and still raises overall
recall by 16\%.

\paragraph{What clustering adds.}
On the reader graph alone, our most important finding is also the least
flattering one. Restricting candidates to a reader's communities gives exactly
the same accuracy as recommending from the same backbone without communities,
because 91\% of the backbone's edge weight already lies inside them: Louvain
recovers structure the backbone already encodes. This changes once the
description similarity is added to the graph. Description links are largely new
(only 6\% of them are also reader links), so in the combined graph only 68\% of
the weight falls inside communities, and the restriction becomes a filter. The
community recommender then beats the same graph without communities by a small
but significant 1.3\%, and finds 6\% more long-tail books than reader-only
communities. The two signals complement rather than replace each other:
descriptions alone reach only a third of the reader graph's recall. Beyond
accuracy, the partition is strong (modularity 0.84), agrees with Amazon's genres
far more than chance (NMI 0.29 against 0.06), and reveals reader groups that cut
across genres, such as post-apocalyptic thrillers read alongside prepper
manuals. Personalisation also matters more than membership: ranking a
community's books by centrality, as first planned, did worse than recommending
the most popular books.

\paragraph{New books.}
For books without readers, SBERT similarity combined with popularity gives the
best hit rate (0.482 against 0.458 for popularity alone). The margin grows to
69\% on the strict test that requires at least two shared readers, which suggests
that descriptions capture shared taste and not merely popularity. The gain is
largest for the least-read new books, exactly the ones popularity cannot help,
but descriptions shorter than about 200 characters carry too little information
to be useful. TF-IDF, which only matches shared words, falls well behind SBERT.

\paragraph{Locality-sensitive hashing.}
LSH proved useful as a candidate generator but not as a final answer. MinHash
on titles proposed 10,728 duplicate pairs, yet only 819 survived a check on
title and description, so its candidates had to be verified. Random-hyperplane
LSH cut the number of comparisons by a factor of 6 to 60 while keeping up to
72\% of the exact nearest neighbours, but with 150k candidates an exact search
took only 18.5 seconds as one matrix multiplication. At this scale exact search
is the better choice; LSH would only pay off on a much larger catalogue.

\paragraph{Limitations.}
The data comes from the 2014 SNAP release and was drawn from its 5-core subset,
so every book had at least five reviews before sampling. Our long tail is
therefore the tail of an already filtered catalogue, and the results may
overstate how well any method handles truly obscure books. Reviews are only a
proxy for purchases, and a liked book is defined purely by a rating of at least
four. This is a coarse signal, since 81\% of all ratings are four or five stars,
and the review text itself (143 words on average) was never used to judge
preference.
% If the sentiment part is added, say here how it refines "liked".
The metadata has no author field and assigns a genre to only about one book in
five, so genre agreement is measured on a subset and series were detected
heuristically from titles. Only 70\% of the books have both a title and a
description, so the description graph leaves the rest isolated. Duplicate editions were merged only in the cold-start
part; the graph results treat editions as separate books, which splits their
readers. Finally, every evaluation is offline and uses a single time split
(January to July 2014), and absolute recall stays low, as expected when choosing
ten books out of 48,000.
